\chapter{Chapter02 — Definitions}
%%% generated by the blueprint pipeline from FABL.Chapter02.SocialChoiceFunctions.Definitions
\section{Overview}

Declarations of the Lean module \texttt{FABL.Chapter02.SocialChoiceFunctions.Definitions}.

\section{Declarations}

\begin{definition}[Threshold sign]
\lean{FABL.thresholdSign}
\statementsource{Analysis-of-Boolean-Functions-by-Ryan-ODonnell}{
  pdf-ffbe9c7b3306-p018-b002,
  pdf-ffbe9c7b3306-p113-b003
}
\label{FABL.thresholdSign}
\leanok
\uses{FABL.Sign}
For $t\in\mathbb R$, the threshold sign of $t$ is $1$ when $t\ge 0$ and $-1$ when $t<0$.
\end{definition}

\begin{theorem}[Threshold sign of a nonnegative number]
\lean{FABL.thresholdSign_of_nonneg}
\label{FABL.thresholdSign_of_nonneg}
\leanok
\uses{FABL.thresholdSign}
\difficulty{1}
For every real number $t\ge 0$, its threshold sign is $1$.
\end{theorem}

\begin{theorem}[Threshold sign of a negative number]
\lean{FABL.thresholdSign_of_neg}
\label{FABL.thresholdSign_of_neg}
\leanok
\uses{FABL.thresholdSign}
\difficulty{1}
For every real number $t<0$, its threshold sign is $-1$.
\end{theorem}

\begin{theorem}[Value of a threshold sign]
\lean{FABL.signValue_thresholdSign}
\label{FABL.signValue_thresholdSign}
\leanok
\uses{FABL.signValue, FABL.thresholdSign}
\difficulty{4}
For every real number $t$, the real value of the sign chosen to be $1$ when $t\ge 0$ and
$-1$ when $t<0$ is $1$ when $t\ge 0$ and $-1$ otherwise.
\end{theorem}

\begin{theorem}[Real value of a sign]
\lean{FABL.signValue_eq_ite}
\label{FABL.signValue_eq_ite}
\leanok
\uses{FABL.Sign, FABL.signValue}
\difficulty{4}
For every sign $s\in\{\pm1\}$, its real value is $1$ if $s=1$ and $-1$ otherwise.
\end{theorem}

\begin{theorem}[Injectivity of sign evaluation]
\lean{FABL.signValue_injective}
\label{FABL.signValue_injective}
\leanok
\uses{FABL.signValue}
\difficulty{3}
The map sending each sign $s\in\{\pm1\}$ to its real value is injective.
\end{theorem}

\begin{theorem}[Monotonicity of the sign value]
\lean{FABL.monotone_signValue}
\label{FABL.monotone_signValue}
\leanok
\uses{FABL.signValue}
\difficulty{4}
The map sending each sign $s\in\{-1,1\}$ to its real value is monotone.
\end{theorem}

\begin{theorem}[Lower bound for a sign]
\lean{FABL.neg_one_le_sign}
\label{FABL.neg_one_le_sign}
\leanok
\uses{FABL.Sign}
\difficulty{4}
Every sign $s\in\{-1,1\}$ satisfies $-1\le s$.
\end{theorem}

\begin{theorem}[A sign is at most one]
\lean{FABL.sign_le_one}
\label{FABL.sign_le_one}
\leanok
\uses{FABL.Sign}
\difficulty{4}
Every integer sign $s\in\{-1,1\}$ satisfies $s\le 1$.
\end{theorem}

\begin{theorem}[Sign bounded above by negative one]
\lean{FABL.sign_eq_neg_one_of_le_neg_one}
\label{FABL.sign_eq_neg_one_of_le_neg_one}
\leanok
\uses{FABL.Sign}
\difficulty{4}
If a sign $s\in\{-1,1\}$ satisfies $s\le -1$, then $s=-1$.
\end{theorem}

\begin{theorem}[A sign above one]
\lean{FABL.sign_eq_one_of_one_le}
\label{FABL.sign_eq_one_of_one_le}
\leanok
\uses{FABL.Sign}
\difficulty{4}
If an integer sign $s\in\{-1,1\}$ satisfies $1\le s$, then $s=1$.
\end{theorem}

\begin{theorem}[Monotonicity of the threshold sign]
\lean{FABL.monotone_thresholdSign}
\label{FABL.monotone_thresholdSign}
\leanok
\uses{FABL.thresholdSign}
\difficulty{4}
The map sending a real number $t$ to $1$ when $t\geq0$ and to $-1$ when $t<0$ is
monotone.
\end{theorem}

\begin{theorem}[Negating a threshold sign]
\lean{FABL.thresholdSign_neg}
\label{FABL.thresholdSign_neg}
\leanok
\uses{FABL.thresholdSign, FABL.thresholdSign_of_neg, FABL.thresholdSign_of_nonneg}
\difficulty{4}
For every nonzero real number $t$, the threshold sign of $-t$ is the negative of the
threshold sign of $t$.
\end{theorem}

\begin{definition}[Positive coordinate set]
\lean{FABL.positiveCoordinateSet}
\label{FABL.positiveCoordinateSet}
\leanok
Let $n$ be a natural number and let $x\in\{-1,1\}^n$. The positive coordinate set of $x$
is the set of indices $i$ for which $x_i=1$.
\end{definition}

\begin{definition}[Positive coordinate count]
\lean{FABL.positiveCoordinateCount}
\label{FABL.positiveCoordinateCount}
\leanok
\uses{FABL.positiveCoordinateSet}
Let $n$ be a natural number and let $x\in\{-1,1\}^n$. The positive coordinate count of
$x$ is the number of indices $i$ for which $x_i=1$.
\end{definition}

\begin{theorem}[Positive coordinate count bound]
\lean{FABL.positiveCoordinateCount_le}
\label{FABL.positiveCoordinateCount_le}
\leanok
\uses{FABL.positiveCoordinateCount}
\difficulty{4}
For any sign vector $x\in\{-1,1\}^n$, the number of its coordinates equal to $1$ is
at most $n$.
\end{theorem}

\begin{definition}[Canonical sign vector with a positive prefix]
\lean{FABL.canonicalCountInput}
\label{FABL.canonicalCountInput}
\leanok
Let $n$ and $k$ be natural numbers. Define a sign vector $x\in\{-1,1\}^n$ by setting
$x_i=1$ when $i<k$ and $x_i=-1$ otherwise.
\end{definition}

\begin{theorem}[Count for the canonical sign vector]
\lean{FABL.positiveCoordinateCount_canonicalCountInput}
\label{FABL.positiveCoordinateCount_canonicalCountInput}
\leanok
\uses{FABL.canonicalCountInput, FABL.positiveCoordinateCount, FABL.positiveCoordinateSet}
\difficulty{1}
For natural numbers $n$ and $k$, let $x\in\{-1,1\}^n$ have value $1$ at coordinate
$i$ exactly when $i<k$, and value $-1$ otherwise. The number of coordinates of $x$ equal
to $1$ is $\min(n,k)$.
\end{theorem}

\begin{theorem}[Monotonicity of canonical count inputs]
\lean{FABL.canonicalCountInput_mono}
\label{FABL.canonicalCountInput_mono}
\leanok
\uses{FABL.canonicalCountInput, FABL.neg_one_le_sign}
\difficulty{4}
Let $n$ be a natural number and let $k,l$ be natural numbers with $k\le l$. The
canonical count input for $k$ is coordinatewise at most the canonical count input for
$l$.
\end{theorem}

\begin{theorem}[Sign sum and positive-coordinate count]
\lean{FABL.sum_signValue_eq_two_mul_positiveCoordinateCount_sub}
\label{FABL.sum_signValue_eq_two_mul_positiveCoordinateCount_sub}
\leanok
\uses{FABL.positiveCoordinateCount, FABL.positiveCoordinateSet, FABL.signValue, FABL.signValue_eq_ite}
\difficulty{4}
Let $n$ be a nonnegative integer and let $x\in\{-1,1\}^n$. The sum of the coordinate
signs of $x$ equals twice the number of coordinates equal to $1$, minus $n$.
\end{theorem}

\begin{theorem}[Nonzero sign sum in odd dimension]
\lean{FABL.sum_signValue_ne_zero_of_odd}
\label{FABL.sum_signValue_ne_zero_of_odd}
\leanok
\uses{FABL.positiveCoordinateCount, FABL.signValue, FABL.sum_signValue_eq_two_mul_positiveCoordinateCount_sub}
\difficulty{4}
Let $n$ be an odd nonnegative integer and let $x\in\{-1,1\}^n$. The sum of the
coordinate signs of $x$ is nonzero.
\end{theorem}

\begin{definition}[Majority function]
\lean{FABL.majority}
\statementsource{Analysis-of-Boolean-Functions-by-Ryan-ODonnell}{
  pdf-ffbe9c7b3306-p043-b004,
  pdf-ffbe9c7b3306-p113-b003,
  pdf-ffbe9c7b3306-p118-b007,
  pdf-ffbe9c7b3306-p119-b001
}
\label{FABL.majority}
\leanok
\uses{FABL.BooleanFunction, FABL.signValue, FABL.thresholdSign}
For each $n\in\mathbb N$, the majority function on $\{-1,1\}^n$ sends an input $x$ to
$1$ when the sum of its coordinates is nonnegative, and to $-1$ otherwise.
\end{definition}

\begin{definition}[Characterization of majority functions]
\lean{FABL.IsMajorityFunction}
\statementsource{Analysis-of-Boolean-Functions-by-Ryan-ODonnell}{pdf-ffbe9c7b3306-p043-b004}
\label{FABL.IsMajorityFunction}
\leanok
\uses{FABL.BooleanFunction, FABL.signValue, FABL.thresholdSign}
For each $n\in\mathbb N$, a function $f:\{-1,1\}^n\to\{-1,1\}$ is a majority function
if, for every input $x$, whenever the sum of the coordinates of $x$ is nonzero, $f(x)$
equals $1$ if that sum is positive and $-1$ if it is negative.
\end{definition}

\begin{theorem}[Majority agrees with non-tied sign majority]
\lean{FABL.majority_isMajorityFunction}
\label{FABL.majority_isMajorityFunction}
\leanok
\uses{FABL.IsMajorityFunction, FABL.majority}
\difficulty{1}
For every natural number $n$, majority on $n$ signs agrees with the threshold sign of
the coordinate sum on every sign vector whose coordinate sum is nonzero.
\end{theorem}

\begin{definition}[Conjunction on the Boolean cube]
\lean{FABL.andFunction}
\statementsource{Analysis-of-Boolean-Functions-by-Ryan-ODonnell}{pdf-ffbe9c7b3306-p044-b001}
\label{FABL.andFunction}
\leanok
\uses{FABL.BooleanFunction}
For each $n\in\mathbb N$, define the conjunction function on $\{-1,1\}^n$ to be $-1$
when every coordinate is $-1$, and $1$ otherwise.
\end{definition}

\begin{definition}[Disjunction on the Boolean cube]
\lean{FABL.orFunction}
\statementsource{Analysis-of-Boolean-Functions-by-Ryan-ODonnell}{pdf-ffbe9c7b3306-p044-b001}
\label{FABL.orFunction}
\leanok
\uses{FABL.BooleanFunction}
For each $n\in\mathbb N$, define the disjunction function on $\{-1,1\}^n$ to be $1$ when
every coordinate is $1$, and $-1$ otherwise.
\end{definition}

\begin{definition}[Coordinate projection]
\lean{FABL.dictator}
\statementsource{Analysis-of-Boolean-Functions-by-Ryan-ODonnell}{
  pdf-ffbe9c7b3306-p044-b002,
  pdf-ffbe9c7b3306-p175-b003
}
\label{FABL.dictator}
\leanok
\uses{FABL.BooleanFunction}
For each $n\in\mathbb N$ and coordinate $i\in\{0,\ldots,n-1\}$, the dictator function
sends an input to its $i$th coordinate.
\end{definition}

\begin{theorem}[Dictator and Singleton Monomial]
\lean{FABL.dictator_toReal_eq_monomial_singleton}
\label{FABL.dictator_toReal_eq_monomial_singleton}
\leanok
\uses{FABL.BooleanFunction.toReal, FABL.dictator, FABL.monomial}
\difficulty{1}
Let $n$ be a nonnegative integer, let $i$ be one of its coordinates, and let $x$ be a
sign vector of length $n$. The real-valued version of the dictator on coordinate $i$ at
$x$ equals the monomial indexed by the singleton set $\{i\}$ at $x$.
\end{theorem}

\begin{definition}[Junta depending on few coordinates]
\lean{FABL.IsKJunta}
\statementsource{Analysis-of-Boolean-Functions-by-Ryan-ODonnell}{pdf-ffbe9c7b3306-p044-b003}
\label{FABL.IsKJunta}
\leanok
For $n\in\mathbb N$, a function $f$ from $\{-1,1\}^n$ to a type $\beta$ is a $k$-junta
if it depends only on some set of at most $k$ coordinates; that is, there is a subset
$S$ of the coordinates with $|S|\le k$ such that $f$ depends only on coordinates in $S$.
\end{definition}

\begin{theorem}[Is k junta iff exists factorization]
\lean{FABL.isKJunta_iff_exists_factorization}
\statementsource{Analysis-of-Boolean-Functions-by-Ryan-ODonnell}{pdf-ffbe9c7b3306-p044-b003}
\label{FABL.isKJunta_iff_exists_factorization}
\leanok
\uses{FABL.IsKJunta, FABL.Sign}
\difficulty{4}
% dataset-pipeline: stub -- prose pending build_informal_statements.py
Statement pending regeneration.
\end{theorem}

\begin{theorem}[A Dictator Is a One-Junta]
\lean{FABL.dictator_isKJunta}
\label{FABL.dictator_isKJunta}
\leanok
\uses{FABL.IsKJunta, FABL.dictator}
\difficulty{3}
Let $n$ be a nonnegative integer and let $i$ be one of its coordinates. The function
that returns the sign of coordinate $i$ is a $1$-junta: it depends only on some set of
at most one coordinate.
\end{theorem}

\begin{definition}[Linear threshold functions]
\lean{FABL.IsLinearThreshold}
\statementsource{Analysis-of-Boolean-Functions-by-Ryan-ODonnell}{
  pdf-ffbe9c7b3306-p044-b005,
  pdf-ffbe9c7b3306-p113-b003,
  pdf-ffbe9c7b3306-p241-b003
}
\label{FABL.IsLinearThreshold}
\leanok
\uses{FABL.BooleanFunction, FABL.signValue, FABL.thresholdSign}
A Boolean function $f\colon\{-1,1\}^n\to\{-1,1\}$ is a linear threshold function if
there are real weights $a_1,\ldots,a_n$ and a real threshold offset $a_0$ such that, for
every input $x$, its value is the sign of $a_0+\sum_{i=1}^n a_i x_i$, with zero assigned
the value $1$.
\end{definition}

\begin{definition}[Equal-weight threshold functions]
\lean{FABL.IsEqualWeightThreshold}
\statementsource{Analysis-of-Boolean-Functions-by-Ryan-ODonnell}{
  pdf-ffbe9c7b3306-p061-b007,
  pdf-ffbe9c7b3306-p113-b003,
  pdf-ffbe9c7b3306-p118-b007
}
\label{FABL.IsEqualWeightThreshold}
\leanok
\uses{FABL.BooleanFunction, FABL.signValue, FABL.thresholdSign}
A Boolean function $f\colon\{-1,1\}^n\to\{-1,1\}$ is an equal-weight threshold function
if there is a real threshold offset $a_0$ such that, for every input $x$, its value is
the sign of $a_0+\sum_{i=1}^n x_i$, with zero assigned the value $1$.
\end{definition}

\begin{theorem}[Majority is a linear threshold function]
\lean{FABL.majority_isLinearThreshold}
\label{FABL.majority_isLinearThreshold}
\leanok
\uses{FABL.IsLinearThreshold, FABL.majority}
\difficulty{3}
For every natural number $n$, majority on $n$ signs is a linear threshold function:
there are a real bias and real coordinate weights such that, on every sign vector,
majority equals the threshold sign of the bias plus the weighted sum of the coordinates.
\end{theorem}

\begin{theorem}[A Dictator Is a Linear Threshold]
\lean{FABL.dictator_isLinearThreshold}
\label{FABL.dictator_isLinearThreshold}
\leanok
\uses{FABL.IsLinearThreshold, FABL.dictator, FABL.signValue, FABL.thresholdSign}
\difficulty{4}
Let $n$ be a nonnegative integer and let $i$ be one of its coordinates. The function
that returns the sign of coordinate $i$ is a linear threshold function: there are a real
bias and real coordinate weights such that, at every sign vector, the function equals
$1$ when the resulting affine linear form is nonnegative and equals $-1$ otherwise.
\end{theorem}

\begin{definition}[Block of a sign vector]
\lean{FABL.inputBlock}
\label{FABL.inputBlock}
\leanok
Let $s$ and $w$ be natural numbers, let $x\in\{-1,1\}^{sw}$, and let
$i\in\{0,\ldots,s-1\}$. The $i$th block of $x$ is the sign vector of length $w$ whose
coordinate $j$ is the coordinate of $x$ indexed by the pair $(i,j)$ under the standard
identification of the $sw$ coordinates with pairs.
\end{definition}

\begin{definition}[Recursive block of a sign vector]
\lean{FABL.recursiveInputBlock}
\label{FABL.recursiveInputBlock}
\leanok
Let $n$ and $d$ be natural numbers, let $x\in\{-1,1\}^{n^{d+1}}$, and let
$i\in\{0,\ldots,n-1\}$. The $i$th recursive block of $x$ is the sign vector of length
$n^d$ obtained by fixing the first index to $i$ under the standard identification of the
coordinates with an index in $\{0,\ldots,n-1\}$ and an index in $\{0,\ldots,n^d-1\}$.
\end{definition}

\begin{definition}[Recursive majority]
\lean{FABL.recursiveMajority}
\statementsource{Analysis-of-Boolean-Functions-by-Ryan-ODonnell}{
  pdf-ffbe9c7b3306-p044-b007,
  pdf-ffbe9c7b3306-p231-b003,
  pdf-ffbe9c7b3306-p239-b004
}
\label{FABL.recursiveMajority}
\leanok
\uses{FABL.BooleanFunction, FABL.majority, FABL.recursiveInputBlock}
For natural numbers $n$ and $d$, define a Boolean function on sign vectors of length
$n^d$ recursively. At depth zero it returns the sole coordinate of the input; at depth
$d+1$, it divides the input into $n$ recursive blocks of length $n^d$, applies the
depth-$d$ function to each block, and returns the majority sign of those $n$ values,
assigning a tie the value $1$.
\end{definition}

\begin{theorem}[Recursive majority at depth zero]
\lean{FABL.recursiveMajority_zero}
\label{FABL.recursiveMajority_zero}
\leanok
\uses{FABL.recursiveMajority}
\difficulty{1}
For every natural number $n$ and every sign vector $x$ with $n^0$ coordinates, the
depth-zero recursive majority is the value of $x$ at its sole coordinate.
\end{theorem}

\begin{theorem}[Recursive majority step]
\lean{FABL.recursiveMajority_succ}
\label{FABL.recursiveMajority_succ}
\leanok
\uses{FABL.majority, FABL.recursiveInputBlock, FABL.recursiveMajority}
\difficulty{1}
For natural numbers $n,d$ and a sign vector $x$ with $n^{d+1}$ coordinates, the
recursive majority of depth $d+1$ is the majority of the $n$ recursive-majority outputs
obtained from the $n$ consecutive blocks of $n^d$ coordinates of $x$.
\end{theorem}

\begin{theorem}[Depth-one recursive majority]
\lean{FABL.recursiveMajority_one}
\label{FABL.recursiveMajority_one}
\leanok
\uses{FABL.majority, FABL.recursiveInputBlock, FABL.recursiveMajority, FABL.recursiveMajority_succ, FABL.recursiveMajority_zero}
\difficulty{3}
For every natural number $n$ and every sign input $x$ indexed by $n^1$ coordinates,
recursive majority of depth one on $x$ equals ordinary $n$-input majority applied to the
signs of $x$ at the corresponding coordinates.
\end{theorem}

\begin{definition}[Tribes function]
\lean{FABL.tribes}
\statementsource{Analysis-of-Boolean-Functions-by-Ryan-ODonnell}{
  pdf-ffbe9c7b3306-p045-b001,
  pdf-ffbe9c7b3306-p096-b005
}
\label{FABL.tribes}
\leanok
\uses{FABL.BooleanFunction, FABL.andFunction, FABL.inputBlock, FABL.orFunction}
For natural numbers $w$ and $s$, the tribes function on sign vectors of length $sw$ is
the disjunction of $s$ conjunctions, each formed from one block of $w$ input
coordinates.
\end{definition}

\begin{theorem}[Count under sign negation]
\lean{FABL.positiveCoordinateCount_negateInput}
\label{FABL.positiveCoordinateCount_negateInput}
\leanok
\uses{FABL.positiveCoordinateCount, FABL.positiveCoordinateCount_le, FABL.signValue, FABL.sum_signValue_eq_two_mul_positiveCoordinateCount_sub}
\difficulty{4}
For any sign vector $x\in\{-1,1\}^n$, the number of coordinates equal to $1$ in the
coordinatewise negation of $x$ equals $n$ minus the number of coordinates equal to $1$
in $x$.
\end{theorem}

\begin{definition}[Permuted input]
\lean{FABL.permuteInput}
\statementsource{Analysis-of-Boolean-Functions-by-Ryan-ODonnell}{pdf-ffbe9c7b3306-p039-b002}
\label{FABL.permuteInput}
\leanok
Given a permutation $\pi$ of the $n$ coordinate indices and a sign vector
$x\in\{-1,1\}^n$, the permuted input is the vector whose coordinate at index $i$ is
$x_{\pi(i)}$.
\end{definition}

\begin{definition}[Permutation of sign vectors]
\lean{FABL.permuteInputEquiv}
\label{FABL.permuteInputEquiv}
\leanok
\uses{FABL.permuteInput}
Let $n$ be a natural number and let $\pi$ be a permutation of $\{0,\ldots,n-1\}$.
Permuting the coordinates of a sign vector by $\pi$ defines a bijection from
$\{-1,1\}^n$ to itself, with inverse given by permuting coordinates by $\pi^{-1}$.
\end{definition}

\begin{definition}[Permutation of a finite set]
\lean{FABL.permuteFinset}
\label{FABL.permuteFinset}
\leanok
For a permutation $\pi$ of the coordinates $\{0,\ldots,n-1\}$ and a subset $S$ of those
coordinates, the permuted subset is $\{\pi(i): i\in S\}$.
\end{definition}

\begin{theorem}[Monomial under coordinate permutation]
\lean{FABL.monomial_permuteInput}
\label{FABL.monomial_permuteInput}
\leanok
\uses{FABL.monomial, FABL.permuteFinset, FABL.permuteInput}
\difficulty{1}
Let $n$ be a natural number, let $\pi$ be a permutation of the $n$ coordinate indices,
let $S$ be a subset of those indices, and let $x$ be a vector of $n$ signs. The monomial
indexed by $S$ at the vector whose $i$th coordinate is $x_{\pi(i)}$ equals the monomial
indexed by the image $\pi(S)$ at $x$.
\end{theorem}

\begin{theorem}[Inverse permutation of a finite set]
\lean{FABL.permuteFinset_symm}
\label{FABL.permuteFinset_symm}
\leanok
\uses{FABL.permuteFinset}
\difficulty{1}
Let $\pi$ be a permutation of $\mathrm{Fin}\,n$, and let $S$ be a finite subset of $\mathrm{Fin}\,n$.
Applying the inverse permutation to $S$ and then applying $\pi$ returns $S$.
\end{theorem}

\begin{theorem}[Fourier coefficients under coordinate permutation]
\lean{FABL.fourierCoeff_comp_permuteInput}
\label{FABL.fourierCoeff_comp_permuteInput}
\leanok
\uses{FABL.fourierCoeff, FABL.monomial, FABL.monomial_permuteInput, FABL.permuteFinset, FABL.permuteFinset_symm, FABL.permuteInput, FABL.permuteInputEquiv}
\difficulty{3}
For every natural number $n$, every permutation $\pi$ of the $n$ coordinates, every
real-valued function $f$ on the sign cube, and every coordinate set $S$, the Fourier
coefficient at $S$ of the function obtained by permuting the input coordinates according
to $\pi$ equals the Fourier coefficient of $f$ at the image of $S$ under $\pi^{-1}$.
\end{theorem}
